\scanpage{59}
% archive sheet number: 58

Par une petite modification de ce processus cela montrera alors que non seulement $\Pi_1(\mathcal X)\longrightarrow\pm1$, mais aussi $\Pi_1(\widetilde{\mathcal X})\longrightarrow\pm1$ est surjectif ; or donc dans le cas $n$ pair, cela implique que $\Pi_1(\unsure{\widetilde{\mathcal X}|\mathcal X|\widetilde{\widetilde{\mathcal X}}|\overset{\approx}{\mathcal X}})\longrightarrow\mathbb D_{2n}$ est surjectif (mais peut-être pas $\Pi_1(\unsure{\overline{\mathcal X}|\widetilde{\mathcal X}})\longrightarrow\mathbb D_{2n}$ ?)

Calcul de la parité du nb d'opérations élémentaires pour passer de $(D_i^{s\unsure{,|\prime|{}}},\omega)$ à $(D_i^{s},-\omega)$.

\figure{\begin{tikzpicture}[x=0.016cm,y=-0.016cm,line cap=round,line join=round]
\draw[thin] (0,37)--(87,37);
\draw[thin] (114,37)--(288,38);
\fill[black] (101,37) circle [radius=1.6pt];
\draw[->] (166,23)--(160,-10);
\node at (100,61) {$s$};
\end{tikzpicture}}

\deleted{Le nb d'opérations est égal à
\[ \begin{aligned}
&2\,(\text{nb de sommets à balayer})\\
&+2\,(\text{nb de droites par dessus lesquelles balayer})\\
&+1\,(\text{décollage})+1\,(\text{atterrissage})
\end{aligned} \]
merveilleux, c'est toujours pair !}

\margin{NB chaque droite par dessus laquelle on balaye compte comme un \uline{autre} sommet par lequel on balayerait}

\deleted{En fait, par transport parallèle, on trouve (sauf erreur \unsure{prov.|p.r.}}

A vrai dire, le raisonnement est \unsure{correct|canulé}, \deleted{par} ce que il faut régler ici dans la version de $\mathcal X$ qui n'a pas de faces supersingulières (pas de ``bandes supersingulières'') car celles-ci sont d'indices $2(n-1)$, non $2n$. Donc il vaut mieux partir d'une position proche de $D_i$ \unsure{passant|pass.nt} par $s$ \deleted{par} un petit pivotement autour de $s$.

Maintenant, \deleted{toute} droite à franchir compte comme \uline{une seule} opération, hélas, et chaque droite \inserted{passant par $s$} \deleted{franchi} est franchie une fois et une seule. Donc le nb correct d'opérations est

\[ \begin{aligned}
&2\,(\text{nb de sommets non alignés sur }s)\\
&+\text{ordre de }s
\end{aligned} \]

donc l'image du circuit dans $\pm1$ est $-1$ ou $+1$, suivant que l'ordre \deleted{est} $\nu(s)$ de $s$ est pair, ou impair. Si donc il y a au moins un sommet simple, on voit que l'on a bien $\Pi_1(\unsure{\mathcal X|\widetilde{\mathcal X}|\widetilde{\widetilde{\mathcal X}}})\longrightarrow\pm1$ surjectif, et aussi $\Pi_1(\widetilde{\mathcal X})\longrightarrow\pm1$ surj.
